\chapter{Methodology}\label{methodology}

\section{Hardware specification}

\begin{itemize}
    \item GPU: 1 $\times$ NVIDIA GeForce RTX 2080 Ti (11 GB VRAM)
    \item CPU: 4 Cores (allocated via SLURM)
    \item Memory: 31 GB RAM
    \item Partition: csedu (Radboud University Science Cluster) 
\end{itemize}

\section{Methodology Super Resolution}
\label{sec:methodology}






The methodology for this project is organized into four main components: the \textbf{Dataset and Preprocessing}, the \textbf{Adaptive-Depth U-Net Architecture}, the \textbf{Experimental Design}, and the \textbf{Training and Evaluation Protocol}. 

The core objective is to investigate the optimal balance between the super-resolution downscale factor ($s$) and the architectural depth of the network.

\subsection{Dataset and Preprocessing}
\label{subsec:data}

The \textbf{DIV2K} dataset, a standard benchmark for single-image super-resolution (SR), was used for training and evaluation. Specifically, the \textbf{Train Data Track 1 – Bicubic Downscaling $\times 4$} subset was utilized, consisting of 800 high-resolution (HR) and low-resolution paired (LR) images for training, and 100 pairs for validation. HR images are approximately $2040 \times 1080$ pixels before cropping, and the LR inputs are $510 \times 270$ pixels (downsampled by a factor of 4).

\subsubsection{Low-Resolution (LR) Data Generation}
\label{subsubsec:lr_generation}

Two primary strategies were considered for generating the low-resolution (LR) training inputs:

\begin{description}
    \item[Synthetic (On-the-fly) Generation] 
    This approach loads only the DIV2K High-Resolution (HR) frames. The \texttt{tf.data} pipeline then crops random $256 \times 256$ patches and immediately applies a bicubic downsample/upsample operation to each patch. This strategy keeps the workflow simple, as no extra assets are required. To maintain consistency, I set the LR degradation to $0.5x$ by default for comparability across scales. The LR inputs are therefore generated dynamically while remaining perfectly pixel-aligned with the HR patches.

    \item[Precomputed (Official) Tiles] 
    This alternative involves using the "official" bicubic-downsampled LR sets provided by DIV2K (e.g., \texttt{X2}, \texttt{X3}, \texttt{X4}). This would require downloading the corresponding directories and pairing every HR crop with its stored LR counterpart. While the benefit is matching the benchmark's exact degradation pipeline, this method is limited to the provided scales (e.g., 0.5 for $\times 2$, $\approx 0.33$ for $\times 3$, 0.25 for $\times 4$) and incurs extra I/O and data management overhead.
\end{description}

For this project, the \textbf{Synthetic (On-the-fly) Generation} method was adopted. All models are trained using LR patches dynamically generated from the HR images, eliminating the need for additional LR datasets.


\subsection{U-Net Architecture}
\label{subsec:architecture}

\paragraph{Vanillin UNet model}

This is the classic 4-level encoder-decoder with symmetric skip connections. Each encoder block consists of a Conv2D $\rightarrow$ BatchNorm $\rightarrow$ Relu and then halves the spatial resolution with MaxPool2D. Each decoder block consists of UpSampling2D (with bilinear interpolation) $\rightarrow$ Conv2D $\rightarrow$  concatenation of the corresponding skip connections. The head is a single 1×1  convolution (with sigmoid activation) that directly predicts the HR RGB patch. 

\paragraph{Modified U-Net with a Residual Head}

The selected architecture is a modified \textbf{U-Net} \cite{ronneberger2015unet}, chosen after reviewing prior work. 
It provides a strong baseline for testing both super-resolution quality and, more importantly, the effectiveness of the proposed adaptive-depth architecture. 
To enhance reconstruction stability, a \textbf{Residual Reconstruction Head} inspired by residual learning \cite{he2016resnet} is applied at the output. 
The final reconstruction $\hat{Y}$ is computed as:
$$
\hat{Y} = \text{ClippedResidualSum}(X_{\text{LR}}, \text{Residual}(X_{\text{LR}}))
$$
where $X_{\text{LR}}$ denotes the low-resolution input and $\text{Residual}(\cdot)$ represents the learned correction. 
The \texttt{ClippedResidualSum} layer ensures that all pixel intensities remain within valid image bounds after residual addition. 

To support arbitrary scaling during both upsampling and downsampling, two custom TensorFlow layers were implemented: 
\texttt{ResizeByScale} and \texttt{ResizeToMatch}. 
Unlike fixed-stride pooling or transposed convolutions, these layers perform continuous spatial scaling using differentiable bilinear interpolation, 
allowing the model to handle non-integer scale factors while maintaining smooth gradient flow.

\paragraph{Adaptive-Depth Principle}

The main idea behind the adaptive-depth design is that the network depth should change with the input scale factor $s$. 
When the input is heavily downscaled, a deep encoder–decoder quickly reduces feature maps to very small sizes, leading to loss of spatial detail. 
For such cases, a shallower network is more efficient and stable. 
At higher scales, more spatial information is available, allowing deeper models to extract richer features. 
However, this increased spatial resolution also raises computational and memory requirements, often resulting in longer training times or the need for smaller batch sizes. 
Based on this intuition, the U-Net depth is adjusted as follows:
\begin{itemize}
    \item \textbf{Small scales ($s \le 0.4$):} Use fewer encoder–decoder levels to preserve detail and reduce computation.  
    \item \textbf{Large scales ($s \ge 0.5$):} Use more levels to capture richer features, at the cost of higher training time and memory use.
\end{itemize}

\subsection{Experimental Design}
\label{subsec:exp_design}

Two primary experiments were designed to isolate and analyze the effects of scale factor and network depth. All models were trained using the \textbf{Charbonnier Loss} (a robust form of L1 loss) for stable optimization.

\subsubsection{Experiment 1: Fixed Depth, Varying Scale}
This experiment establishes a baseline performance across a wide range of scale factors ($s$) using a uniform network capacity.
\begin{itemize}
    \item \textbf{Network Depth:} Fixed at \textbf{3 encoder-decoder levels} for all runs.
    \item \textbf{Scale Factors ($s$):} Varied across a broad spectrum: $\{0.2, 0.3, 0.4, 0.5, 0.6, 0.7, 0.8, 0.9\}$.
\end{itemize}
The results from this experiment reveal the inherent performance ceiling of a fixed architecture when faced with increasingly higher-fidelity inputs.

\subsubsection{Experiment 2: Adaptive Depth per Scale}
This is the core test of the proposed adaptive architecture. The network depth is tuned for each scale factor based on the principle described in Section \ref{subsec:architecture}, aiming to optimize feature map size stability and capacity.

\begin{table}[h!]
    \centering
    \caption{Adaptive-Depth Configuration per Scale Factor}
    \label{tab:adaptive_depth}
    \begin{tabular}{|c|c|c|}
        \hline
        \textbf{Scale ($s$)} & \textbf{Depth (Levels)} & \textbf{Rationale} \\
        \hline
        0.20 & 1 & Maintain feature maps $\ge 50$ px \\
        0.30 & 2 & Adds capacity; avoids map collapse \\
        0.40 & 3 & Balanced configuration \\
        0.50 & 3 & Baseline stability \\
        0.60 & 4 & Increased capacity for higher detail \\
        0.70 & 5 & Significant capacity increase \\
        0.80 & 5 & High-detail retention \\
        \hline
    \end{tabular}
\end{table}

The performance of these adaptive models is benchmarked against the fixed-depth models (Experiment 1) and a \textbf{Bicubic Baseline} (i.e., the LR image upsampled by bicubic interpolation to HR size).


\subsubsection{Evaluation of the models} 

The evaluation was conducted using the Test Data Track 1 bicubic downscaling (x2) subset. Due to limited access to the full dataset, these images served as the Low-Resolution (LR) inputs. They were subsequently downsampled further using the Synthetic (On-the-fly) Generation strategy described in Section \ref{subsec:data}.

\section{Methodology: Segmentation}
\label{sec:methodology_segmentation}

\subsection{Dataset and Preprocessing}

\paragraph{Dataset}
ISIC-2017 lesion segmentation dataset is a common dataset used in research into segmentation and U-Net design. The selection of this dataset was done as it was easily comparable to the reearch in the related work section. The dataset consists of 3 splits train (2000 images), test (600 images), and valid (150 images). 

\paragraph{Preprocessing}
All images and masks undergo the following standardization before entering the network

Input Images:
\begin{itemize}
    \item Resizing: Images are resized to a fixed input size of 256 $\times$ 256 pixels using Area interpolation (which is preferred for downsampling to avoid aliasing artifacts).
    \item Normalization: Pixel values are normalized to the range [0, 1] (converted to float32).
    \item Channels: Loaded as 3-channel RGB images.
\end{itemize}



\paragraph{Segmentation Masks}

\begin{itemize}
    \item Resizing: Resized to 256 $\times$ 256 pixels using Nearest Neighbor interpolation. This is crucial to preserve the discrete nature of the class labels and avoid introducing intermediate values at boundaries.
    \item Binarization: The masks are thresholded at 0.5 to ensure strict binary values ({0.0, 1.0}), representing background and lesion respectively. 
\end{itemize}

\paragraph{Data Augmentation}

To prevent overfitting, the following augmentations are applied dynamically during training (and disabled during validation/testing):

\begin{itemize}
    \item Random Rotation: The image and mask are rotated by 0, 90, 180, or 270 degrees ($k \in {0, 1, 2, 3}$).
    \item Random Flips: Horizontal Flip applied with $50\%$ probability and Vertical Flip applied with $50\%$ probability.
    \item Random Scale qne Crop:The image is scaled up by a random factor between 1.0x and 1.15x. It is then randomly cropped back to the original size (256 $\times$ 256).
    \item Note: Bilinear interpolation is used for scaling images, while Nearest Neighbor is used for masks.
\end{itemize}

\paragraph{Pipeline Optimization}

The data is loaded using the tf.data API with the following performance optimizations:
\begin{itemize}
    \item Shuffling: The training dataset is fully shuffled at each epoch.
    Parallel Mapping: Preprocessing and augmentation are parallelized $(num\_parallel\_calls=tf.data.AUTOTUNE)$
    \item Prefetching: Data is prefetched to the GPU memory to prevent I/O blocking.
\end{itemize}



\subsection{U-Net Architecture}

\paragraph{Vanilla U-Net}

\begin{enumerate}
    \item \textbf{Core Structure}: The model follows the classic Encoder-Decoder (U-Net) architecture with skip connections.
    \begin{itemize}
        \item Input: RGB Images of shape $(256, 256, 3)$.
        \item Output: Binary segmentation mask of shape $(256, 256, 1)$.
    \end{itemize}

    \item \textbf{Building Blocks}: The architecture is built using a standard Convolutional Block (\texttt{conv\_block}) repeated throughout the network.
    \begin{itemize}
        \item \textbf{Double Convolution}: Each block consists of two consecutive sets of:
        \begin{itemize}
            \item \texttt{Conv2D}: $3 \times 3$ kernel, \texttt{padding="same"} (preserves spatial dimensions).
            \item \texttt{LayerNormalization}: Applied over the channels (\texttt{axis=-1}).
            \item \texttt{ReLU Activation}: Non-linear activation function.
        \end{itemize}
    \end{itemize}

    \item \textbf{Encoder (Contracting Path)}: The encoder progressively reduces spatial resolution while increasing feature depth.
    \begin{itemize}
        \item \textbf{Depth}: Configurable, default is 4 levels.
        \item \textbf{Downsampling}: Uses \texttt{MaxPooling2D} with a $2 \times 2$ pool size.
        \item \textbf{Feature Channels}: Starts with \texttt{base\_channels} (default 32) and doubles at each level ($32 \rightarrow 64 \rightarrow 128 \rightarrow 256$).
        \item \textbf{Skip Connections}: The output of each encoder block (before pooling) is saved to be concatenated with the corresponding decoder level.
    \end{itemize}

    \item \textbf{Bottleneck}: A single \texttt{conv\_block} at the bottom of the "U".
    \begin{itemize}
        \item Contains the highest number of filters (e.g., 512 if depth is 4).
    \end{itemize}

    \item \textbf{Decoder (Expansive Path)}: The decoder recovers spatial resolution to generate the segmentation mask.
    \begin{itemize}
        \item \textbf{Upsampling}: Uses \texttt{Conv2DTranspose} (Deconvolution) with a $2 \times 2$ kernel and stride 2.
        \item \textbf{Concatenation}: The upsampled feature map is concatenated with the corresponding skip connection from the encoder.
        \item \textbf{Processing}: Followed by a \texttt{conv\_block} to process the combined features.
        \item \textbf{Feature Channels}: Halves at each level as it goes up.
    \end{itemize}

    \item \textbf{Output Layer}
    \begin{itemize}
        \item \textbf{Convolution}: A final \texttt{Conv2D} with a $1 \times 1$ kernel.
        \item \textbf{Activation}: Sigmoid activation to squash values between $[0, 1]$ for binary classification (pixel-wise probability).
    \end{itemize}
\end{enumerate}


\paragraph{Adaptive-Depth U-Net Architecture}


\begin{enumerate}
    \item \textbf{Normalization Strategy}
    \begin{itemize}
        \item \textbf{Adaptive U-Net}: Uses \texttt{BatchNormalization}. This normalizes across the entire batch, which is generally standard for CNNs but sensitive to batch size.
        \item \textbf{Vanilla}: Used \texttt{LayerNormalization} (normalizes per sample).
    \end{itemize}

    \item \textbf{Upsampling Method}
    \begin{itemize}
        \item \textbf{Adaptive U-Net}: Uses Bilinear Upsampling (\texttt{UpSampling2D}). This is a fixed, non-learnable interpolation operation. It is often smoother and avoids the "checkerboard artifacts" sometimes caused by learnable deconvolutions.
        \item \textbf{Vanilla}: Used \texttt{Conv2DTranspose} (learnable deconvolution).
    \end{itemize}

    \item \textbf{Regularization (Dropout)}
    \begin{itemize}
        \item \textbf{Adaptive U-Net}: Includes \texttt{SpatialDropout2D}. It drops entire feature maps (channels) rather than individual pixels, which is more effective for convolutional feature maps where adjacent pixels are highly correlated.
        \begin{itemize}
            \item Applied after the bottleneck.
            \item Applied in the decoder blocks after concatenation.
        \end{itemize}
        \item \textbf{Vanilla}: Had no dropout.
    \end{itemize}

    \item \textbf{Depth Flexibility}
    \begin{itemize}
        \item \textbf{Adaptive U-Net}: While both scripts accept a \texttt{depth} argument, this script explicitly names the model \texttt{adaptive\_unet\_depth\{N\}}, highlighting that the architecture is designed to be dynamically scaled (e.g., changing depth to 3, 4, or 5) as a primary hyperparameter for the experiments.
    \end{itemize}
\end{enumerate}

\subsection{Running Optuna to Obtain the Best Hyperparameter}
\label{subsec:optuna_hyperparams}

The segmentation model’s hyperparameter were optimized using the \textbf{Optuna} framework, which performed a Bayesian search over a predefined parameter space.  
The best-performing configuration was obtained after multiple trials and provides us with a good base for training your Segmentation model, yielding the following parameters:

\begin{itemize}
    \item \textbf{Learning rate:} $3.0592 \times 10^{-5}$
    \item \textbf{Base channels:} 32
    \item \textbf{Depth:} 5
    \item \textbf{Batch size:} 8
    \item \textbf{Data augmentation:} True
\end{itemize}



\subsection{Experimental Design}
\label{subsec:exp_design}

Similar approach to SR was implemented with the loss function being Dice and IoU. 



\subsubsection{Experiment 1: Fixed Depth, Varying Scale} This experiment establishes a baseline performance across a wide range of scale factors ($s$) using a uniform network capacity. The network depth was fixed at 4 encoder-decoder levels, resulting in approximately 31 million trainable parameters.

\begin{itemize} 
    \item \textbf{Architecture:} Standard U-Net with 4 levels (Base Channels: 64). 
    \item \textbf{Scale Factors ($s$):} Varied across ${0.2, 0.3, \dots, 0.9}$. 
    \item \textbf{Training Protocol:} 
    \begin{itemize} 
        \item \textbf{Optimizer:} AdamW ($\lambda=1e-4$) with Cosine Decay learning rate schedule (initial $lr=1e-3$). 
        \item \textbf{Loss Function:} Hybrid Loss ($0.4 \cdot \text{CE} + 0.6 \cdot \text{Dice}$). 
        \item \textbf{Image Resolution:} $256 \times 256$ pixels (ISIC 2017 Dataset). 
        \item \textbf{Batch Size:} Adjusted between 1 and 8 depending on scale factor to fit within GPU memory (NVIDIA 2080 Ti). \item \textbf{Duration:} Trained for up to 100 epochs with early stopping (patience=15). 
    \end{itemize} 
\end{itemize}


\subsubsection{Experiment 2: Adaptive Depth per Scale} This is the core test of the proposed adaptive architecture. The network depth is tuned for each scale factor to optimize feature map size stability. The goal is to maintain the bottleneck spatial resolution roughly between $5 \times 5$ and $25 \times 25$ pixels, ensuring sufficient receptive field without collapsing spatial dimensions. 

\begin{itemize} 
    \item \textbf{Architecture:} Adaptive Depth U-Net (Depth varies from 1 to 5 levels; Base Channels: 64). 
    \item \textbf{Scale Factors ($s$):} Varied across ${0.2, 0.3, \dots, 0.8}$. 
    \item \textbf{Training Protocol:} 
        \begin{itemize} 
            \item \textbf{Optimizer:} AdamW ($\lambda=1e-4$) with Cosine Decay learning rate schedule (initial $lr=1e-3$). \item \textbf{Loss Function:} Hybrid Loss ($0.4 \cdot \text{CE} + 0.6 \cdot \text{Dice}$). 
            \item \textbf{Image Resolution:} $256 \times 256$ pixels (ISIC 2017 Dataset). 
            \item \textbf{Batch Size:} Adjusted between 1 and 8 depending on scale factor to fit within GPU memory (NVIDIA 2080 Ti). 
            \item \textbf{Duration:} Trained for up to 100 epochs with early stopping (patience=30).
            \item \textbf{Note:} The adative model were taking much longer to train as the scale increased so i had to increase the patience 
        \end{itemize} 
\end{itemize}

\begin{table}[h!] 
    \centering 
    \caption{Adaptive-Depth Configuration per Scale Factor} \label{tab:adaptive_depth} 
    \begin{tabular}{|c|c|c|} 
    \hline 
    \textbf{Scale ($s$)} & \textbf{Depth (Levels)} & \textbf{Rationale (Bottleneck Size)} \\
    \hline 
    0.20 & 1 & $51 \to 25$ px (Avoids collapse to $3$ px) \\ 
    0.30 & 2 & $77 \to 19$ px \\ 
    0.40 & 3 & $102 \to 12$ px \\
    0.50 & 3 & $128 \to 16$ px \\
    0.60 & 4 & $154 \to 9$ px \\ 
    0.70 & 5 & $179 \to 5$ px \\ 
    0.80 & 5 & $205 \to 6$ px \\ 
    \hline 
    \end{tabular} 
\end{table}



\subsubsection{Evaluation of the models}

Evaluation of the model was stair forward was we have 600 test images in which we used both IoU and dice as a loss measure. 